\section{A global character from the two elliptic frames}
\label{sec:global-frame-character}

We define the frame signs, extend their sum to every multiplier, and
fix its arithmetic normalization. The construction uses the normalized
frame vanishing theorem of \cite[Theorem~3]{H17}. The extension and
the conjugation invariance are proved here.

\subsection{The frames and their normalization}

Let $\widehat B_3=\langle u,v\mid uvu=vuv\rangle$ be the profinite
braid group and put $c=(uv)^3$. Its central quotient is
\[
 \Gamma=\widehat B_3/\overline{\langle c\rangle}
      =\langle s,r\mid s^2=r^3=1\rangle,
 \qquad s=uvu,\quad r=uv,\quad \tau=sr=u^{-1}.
\]
The free pure subgroup $F=\ker(\Gamma\to S_3)$ has generators
$x=u^2=\tau^{-2}$ and $y=v^2=sxs$. We also use the subgroup
$\langle u^2,v^2\rangle\subset\widehat B_3$, which maps isomorphically
to $F$. A gentle element $g=(\lambda_g,f_g)$ acts on the braid
group by
\[
 \Phi_g(u)=u^{\lambda_g},\qquad
 \Phi_g(v)=f_g(u^2,v^2)^{-1}v^{\lambda_g}f_g(u^2,v^2),
 \qquad \Phi_g(c)=c^{\lambda_g}.
\]
These are automorphisms and induce the identity on the permutation
quotient \cite[Section 3, before Lemma~5]{H17}.

The finite-subgroup and centralizer theorem for the free profinite
product gives unique pure frames
\[
 \Phi_g(s)=\beta_gs\beta_g^{-1}\quad(g\in\GTz),
 \qquad
 \Phi_h(r)=A_hrA_h^{-1}\quad(\lambda_h=1),
 \qquad \beta_g,A_h\in F.
\]
These are the frames of \cite[Lemma~5]{H17}. For the cubic frame,
the multiplier-one condition fixes the conjugacy class of $r$ in
$\Gamma^{\mathrm{ab}}=C_6$, distinguishing it from $r^{-1}$.
The frames vary continuously: their graphs are closed in compact
products, and uniqueness makes the corresponding projections
homeomorphisms. Their composition identities are
\begin{equation}
 \beta_{gh}=\Phi_g(\beta_h)\beta_g,
 \qquad A_{hk}=\Phi_h(A_k)A_h
       \quad(\lambda_h=\lambda_k=1).
 \label{eq:frame-cocycles}
\end{equation}

Let $\theta$ interchange $x,y$. Comparison of the two expressions
for $\Phi_g(y)$ gives
$f_g=\theta(\beta_g)\beta_g^{-1}$. Indeed their possible discrepancy
centralizes $y^{\lambda_g}$; abelianization makes its exponent zero.
Thus the two coordinates of $[\beta_g^{-1}]\in F^{\mathrm{ab}}$
are equal. Define $\rho_2(g)$ by
\[
 [\beta_g^{-1}]=(\rho_2(g),\rho_2(g)).
\]
It is a continuous cocycle,
\[
 \rho_2(gh)=\rho_2(g)+\lambda_g\rho_2(h).
\]
Its restriction to $G_{\mathbb Q}$ is the Kummer cocycle of the
coherent positive roots of two \cite[Lemma~7]{H17}.
Write
\[
 \mathcal G=\GTz,\qquad
 N_0=\{g\in\mathcal G:\lambda_g=1,\ \rho_2(g)=0\},
 \qquad H=G_{\mathbb Q}.
\]
The cocycle makes $N_0$ closed and normal.

We will use the arithmetic decomposition
\begin{equation}
 \mathcal G=N_0H,\qquad
 N_0\cap H=G_{K_2},\qquad
 K_2=\mathbb Q(\mu_\infty,2^{1/\infty}).
 \label{eq:gentle-arithmetic-normalization}
\end{equation}
This is Proposition~\ref{prop:normalized-comparison}; its intersection formula follows from the two defining Galois actions.

\subsection{The signed matrices at level twelve}

Put
\[
 U=\begin{pmatrix}1&1\\0&1\end{pmatrix},\qquad
 L=\begin{pmatrix}1&0\\-1&1\end{pmatrix},\qquad
 S=ULU,\qquad R_+=UL,\qquad R=-R_+.
\]
The assignments $u\mapsto U$, $v\mapsto L$ define the signed
representation $\widetilde q_\ell:\widehat B_3\to
\operatorname{SL}_2(\mathbb Z_\ell)$ for every prime $\ell$.
Its restriction to $F$ is denoted $q_\ell$. Reducing the signed
representation modulo twelve gives $q_{12}$ on $F$ and a projective
representation $q$ of $\Gamma$.
The signed pure image is
\[
 E_{12}=V_4\times\operatorname{SL}_2(\mathbb F_3),
 \qquad V_4=\{I,U^2,L^2,U^2L^2\}\pmod4.
\]
The matrix $-I$ is not in $V_4$, so this image embeds in the
projective group modulo twelve. Its projective centralizers of
$R_+$ and $S$ are respectively
\begin{equation}
 \{I\}\times\bigl(\langle-I\rangle\times\langle R\rangle\bigr),
 \qquad
 \{I,U^2L^2\}\times\langle S\rangle.
 \label{eq:finite-frame-centralizers}
\end{equation}
These follow directly by multiplying the indicated matrices; they
are the finite groups in \cite[Lemma~12]{H17}.

\begin{lemma}\label{lem:frame-all-multipliers-congruence}
For every $g\in\mathcal G$ and every prime $\ell$ there is a unique
matrix
\[
 D_{g,\ell}=\begin{pmatrix}\lambda_{g,\ell}&b_{g,\ell}\\0&1\end{pmatrix}
\]
such that
$\widetilde q_\ell\Phi_g=\operatorname{Ad}(D_{g,\ell})\widetilde q_\ell$.
It depends continuously on $g$, and $D_{gh,\ell}=D_{g,\ell}D_{h,\ell}$.
For $h\in N_0$, its reduction modulo twelve is the identity.
\end{lemma}

\begin{proof}
Put $a=\lambda_{g,\ell}$ and
$M=\widetilde q_\ell(\Phi_g(v))=
\left(\begin{smallmatrix}\alpha&b\\d&2-\alpha\end{smallmatrix}\right)$.
It is conjugate to $L^a$, so its trace is two and its determinant
is one. The lower-left entry of the braid relation
$U^aMU^a=MU^aM$ gives $d(1+ad)=0$.
If $d=0$, the domain property of $\mathbb Z_\ell$, together with
the trace and determinant, makes $M$ upper unipotent. The images
of both braid generators would commute, contradicting surjectivity
onto $\operatorname{SL}_2(\mathbb Z_\ell)$. Hence $d=-a^{-1}$
and $b=a(\alpha-1)^2$. Direct multiplication gives
\[
 D_{g,\ell}=
  \begin{pmatrix}a&a(1-\alpha)\\0&1\end{pmatrix}.
\]
This proves existence and continuity. A matrix commuting with both
$U,L$ is scalar; the specified lower-right entry gives uniqueness.
Composition then gives the product identity.

For $\lambda_h=1$, write the reduction modulo twelve as $U^{t_h}$.
Modulo four the conjugator $f_h\in F'$ maps trivially because
$V_4$ is abelian, so $t_h=0\pmod4$. Modulo three,
$U^{-t_h}q_3(\beta_h)$ lies in the quaternion subgroup of
$\operatorname{SL}_2(\mathbb F_3)$: its projective class centralizes
the involution represented by $S$. In the quotient of order three,
the classes of $U$ and $L$ agree, whereas
$[\beta_h]=(-\rho_2(h),-\rho_2(h))$ gives class
$-4\rho_2(h)=-\rho_2(h)$. Thus
$t_h=-\rho_2(h)\pmod3$. This is the shear formula of
\cite[Lemma~13]{H17}; in particular $D_h=I\pmod{12}$ on $N_0$.
\end{proof}

Define the closed subgroup
\[
 H_0=\{h\in\mathcal G:\lambda_h=1,\ t_h=0\pmod{12}\}.
\]
It contains $N_0$. For $h\in H_0$, the two frames lie in
\eqref{eq:finite-frame-centralizers}, so define
\begin{equation}
 q_3(A_h)\in(-I)^{a_3(h)}\langle R\rangle,
 \qquad
 q_3(\beta_h)\in S^{a_2(h)}\langle-I\rangle,
 \qquad a_2(h),a_3(h)\in\mathbb F_2.
 \label{eq:frame-sign-definition}
\end{equation}
These are exactly the signs of \cite[Definition~14]{H17}.
The cocycles \eqref{eq:frame-cocycles}, the identity action on
$E_{12}$, and the abelian groups in
\eqref{eq:finite-frame-centralizers} make $a_2,a_3$ continuous
characters on $H_0$. Put $\delta=a_2+a_3$.

The sign convention can also be recorded by a single unit. Let
$\xi(h)\in(\mathbb Z/12\mathbb Z)^\times$ be determined by
\[
 \xi(h)\equiv(-1)^{a_2(h)}\pmod4,\qquad
 \xi(h)\equiv(-1)^{a_3(h)}\pmod3.
\]
For the quadratic Dirichlet character $\chi_{12}$ the complete table is
\[
 \begin{array}{c|rrrr}
  t&1&5&7&11\\ \hline
  \chi_{12}(t)&1&-1&-1&1
 \end{array}
 \qquad (-1)^{\delta(h)}=\chi_{12}(\xi(h)).
\]
All elements of $H_0$ have multiplier one. The unit $\xi$ records
the two frame signs separately from that multiplier.

\subsection{Conjugation by elements of every multiplier}

\begin{proposition}\label{prop:frame-sign-conjugation}
The restrictions of $a_2,a_3$, and hence $\delta$, to $N_0$ are
invariant under conjugation by $\mathcal G$.
\end{proposition}

\begin{proof}
Let $h\in N_0$, $g\in\mathcal G$, and $\eta=ghg^{-1}$.
Normality gives $\eta\in N_0$. All matrix calculations in this
proof are modulo three. Write $\overline D_g$ for the reduction
of $D_{g,3}$.

Applying $\Phi_\eta\Phi_g=\Phi_g\Phi_h$ to $s$ and comparing
unique pure frames gives
\[
 \beta_\eta=\Phi_\eta(\beta_g)^{-1}\Phi_g(\beta_h)\beta_g.
\]
Since $\overline D_\eta=I$,
\[
 q_3(\beta_\eta)=q_3(\beta_g)^{-1}\overline D_g
      q_3(\beta_h)\overline D_g^{-1}q_3(\beta_g).
\]
Thus the traces of their squares agree. The defining groups give
\[
 \operatorname{tr}(q_3(\beta_h)^2)=2(-1)^{a_2(h)},
\]
which proves invariance of $a_2$.

Put $\epsilon\in\{1,-1\}$ equal to $\lambda_g$ modulo three.
The torsion theorem gives
$\Phi_g(r)=\gamma_g r^\epsilon\gamma_g^{-1}$.
The exponent follows from the action on $\Gamma^{\mathrm{ab}}=C_6$.
Since $\Phi_g$ induces the identity on $S_3$, multiplying
$\gamma_g$ by a power of $r$ permits the form
$\gamma_g=B_gs^e$, with $B_g\in F$ and $e=0$ or one.
Here $e=0$ when $\epsilon=1$ and $e=1$ when $\epsilon=-1$.
Comparison of the two actions on $r$ gives
\begin{equation}
 A_\eta=\Phi_\eta(\gamma_g)^{-1}\Phi_g(A_h)\gamma_g.
 \label{eq:cubic-frame-conjugation}
\end{equation}
The right side is pure and conjugates $r$ to $\Phi_\eta(r)$,
so equality follows from uniqueness of the pure frame.

We check this equality in the signed representation. Lift
$\gamma_g$ to $\widetilde\gamma_g=B_g(uvu)^e$ in
$\widehat B_3$. Conjugation by $uvu$ interchanges $x,y$ and
normalizes the embedded free subgroup. Since $\lambda_\eta=1$,
\[
 \Phi_\eta(uvu)=\beta_\eta(uvu)\beta_\eta^{-1}.
\]
Both sides project to the same modular element and have square $c$.
A central discrepancy would have square one, and the central
procyclic group $\overline{\langle c\rangle}\simeq\widehat{\mathbb Z}$
is torsion free. This proves the equality.
It follows that the lifted right side of
\eqref{eq:cubic-frame-conjugation} belongs to the embedded free
subgroup: the two factors involving $uvu$ cancel through its
normalizing action. Its projection equals $A_\eta$, and the free
subgroup injects into the central quotient, so it is the signed
lift of $A_\eta$ itself.

Applying the signed matrix representation therefore gives
\[
 q_3(A_\eta)=
 \widetilde q_3(\widetilde\gamma_g)^{-1}\overline D_g
 q_3(A_h)\overline D_g^{-1}
 \widetilde q_3(\widetilde\gamma_g).
\]
The trace is unchanged. Every power of $R$ has trace two in
$\mathbb F_3$, so
$\operatorname{tr}(q_3(A_h))=2(-1)^{a_3(h)}$.
This proves invariance of $a_3$ as well.
\end{proof}

\begin{theorem}[Arithmetic normalization of the frame character]
\label{thm:global-frame-character}
There is a unique continuous character
\[
 \Delta:\mathcal G\longrightarrow\mathbb F_2
\]
which restricts to $\delta$ on $N_0$ and is zero on $G_{\mathbb Q}$.
For any $g\in\mathcal G$, choose $\sigma\in G_{\mathbb Q}$
with $(\lambda,\rho_2)(\sigma)=(\lambda,\rho_2)(g)$. Then
\[
 \Delta(g)=\delta(g\sigma^{-1}).
\]
This character vanishes on the full pentagon group.
\end{theorem}

\begin{proof}
By \cite[Theorem~3]{H17}, $\delta$ vanishes on the normalized
pentagon subgroup. In particular it vanishes on
$N_0\cap H=G_{K_2}$. Using
\eqref{eq:gentle-arithmetic-normalization}, define
$\Delta(n\sigma)=\delta(n)$ for $n\in N_0$, $\sigma\in H$.
Two decompositions differ by an element of $N_0\cap H$, so this
is well defined. Proposition~\ref{prop:frame-sign-conjugation}
proves its homomorphism law. The multiplication map
$N_0\times H\to\mathcal G$ is a quotient map, since it is a
continuous surjection from a compact space to a Hausdorff space.
Thus the character is continuous. The two specified restrictions
and the decomposition prove uniqueness, as well as the displayed
formula. Finally \cite[Theorem~2]{H17} decomposes the pentagon
group into its normalized subgroup and $G_{\mathbb Q}$, both of
which have zero image. Hence $\Delta$ vanishes on $\GT$.
\end{proof}
